% ============================================================
\section{Motivation and Goals}

Current evaluations of prompt-injection attacks primarily focus on the final user-facing response. 
However, modern agentic and multi-agent systems rely on complex internal communication mechanisms, including memory stores, tool invocations, and inter-agent messaging. 
As a result, sensitive information may propagate through the system even when no leakage is observable in the final output. 
This raises concerns that existing evaluation methodologies may underestimate attack impact and real exposure when they inspect only the final user-visible output. 
Furthermore, manually designing attacks for such systems is labor-intensive and may fail to uncover subtle attack strategies. 
Although prior work has investigated adaptive attacks on inter-agent messages
and has benchmarked leakage in internal traces, these directions have not
been integrated into an AutoResearch-style loop that uses execution feedback
to search a structured attack space. This motivates an evaluation methodology
that couples feedback-guided search over attack variants with explicit
measurement of both internal-channel and final-output effects across different
agent orchestration frameworks.

\paragraph{Research Questions:}
\begin{itemize}[leftmargin=2cm]
    \item[\textbf{RQ1:}] How does internal-channel-aware evaluation change the
            measured leakage risk compared with final-output-only evaluation in
            agentic and multi-agent LLM systems?
    \item[\textbf{RQ2:}] To what extent do adversarial attack variants degrade
            task utility or increase operational burden in agentic and
            multi-agent LLM systems?
    \item[\textbf{RQ3:}] How do leakage, final-output exposure, and degradation
            effects vary across different agentic and multi-agent orchestration
            frameworks?
\end{itemize}
